\documentclass[11pt,a4paper]{article}
\usepackage[hmargin=22mm,vmargin=21mm]{geometry}
\usepackage[T1]{fontenc}
\usepackage{lmodern,amsmath,amssymb,amsthm,mathtools,enumitem,microtype,hyperref,booktabs}
\hypersetup{colorlinks=true,linkcolor=blue,urlcolor=blue,citecolor=blue,pdftitle={Certified Low-Degree Coefficient-Minor Positivity for Fibonacci Matching Transforms},pdfauthor={Ryutaro Yonezu}}
\newtheorem{theorem}{Theorem}
\newtheorem{lemma}[theorem]{Lemma}
\newtheorem{corollary}[theorem]{Corollary}
\newcommand{\coef}[2]{[u^{#1}]#2}
\newcommand{\M}{\mathcal M}
\title{Certified Low-Degree Coefficient-Minor Positivity\\for Fibonacci Matching Transforms\\\large A Computer-Assisted Result for $3\leq k\leq180$}
\author{Ryutaro Yonezu\\\small Independent Researcher}
\date{10 November 2026\\\small Preprint v1.0}
\begin{document}
\maketitle
\begin{abstract}
For a matching polynomial $Z_m(u)=\sum_j\binom{m-j}{j}u^j$ and a binomially
transformed polynomial $p_{n,k}(u)=\sum_j\binom{k+2-j}{j}\binom{n-k-2+j}{k}u^j$,
we study adjacent $2\times2$ coefficient minors of two auxiliary products
$A$ and $B$. An exact, reproducible Newton-basis certificate establishes
strict positivity of all minors indexed $0\le r\le k$, uniformly for every
integer $n\ge2k+4$, for each $k$ in the certified range $3\le k\le180$.
A conditional one-factor reduction converts the necessary sign tests into
smaller polynomial certificates. A separate direct construction of the
original minors verifies $3\le k\le22$ without that reduction.
The motivating coefficient family specializes at $u=1$ to big-descent
enumeration on permutations avoiding 123 and 132. The present result
neither proves the conjectured log-concavity of that distribution nor
establishes positivity for every $k$. Novelty remains unconfirmed.
\end{abstract}
\noindent\textbf{Keywords:} Fibonacci polynomials; matching polynomials;
total positivity; coefficient minors; Newton basis; computer-assisted proof;
big descents; pattern avoidance.\\
\textbf{MSC (provisional):} 05A05, 05A15, 05A20.

\section{Background and precise scope}
Elizalde, Rivera Jr., and Zhuang~\cite{ERZ} derive a bivariate generating
function for big descents of permutations avoiding $\{123,132\}$
(Their Theorem~3.5) and conjecture log-concavity for broad classes
(Their Conjecture~4.6). Here we study a different, coefficientwise
positivity property in the auxiliary matching-polynomial variable $u$.
It is not an assertion about the coefficients of the original big-descent
generating polynomial as a sequence in $k$.
The reason for studying these minors is intrinsic to the transformed
polynomials: positivity of their low-degree minors gives a strict ordering
of successive coefficient ratios and hence monotonicity of a truncated
quotient in the matching variable $u$ (Section~\ref{sec:meaning}).
The example in Section~\ref{sec:crosschecks} shows that, at the boundary
index $r=k$, multiplication by the common matching-polynomial factor
is genuinely needed: the unsmoothed minor $\M_k(P,Q)$ can be negative.

Throughout, binomial coefficients outside their natural integer range
are zero. Define, for $m\ge0$ and relevant $n,k$,
\begin{align}
 Z_m(u)&=\sum_{j=0}^{\lfloor m/2\rfloor}\binom{m-j}{j}u^j,\label{eq:Z}\\
 p_{n,k}(u)&=\sum_{j=0}^{\lfloor(k+2)/2\rfloor}
 \binom{k+2-j}{j}\binom{n-k-2+j}{k}u^j.\label{eq:p}
\end{align}
For fixed $n,k$, set
\begin{align}
 P(u)&=p_{n,k}(u)^2, & Q(u)&=p_{n,k-1}(u)p_{n,k+1}(u),\\
 K(u)&=Z_{k+2}(u)^2, &
 A(u)&=Z_{k+1}(u)Z_{k+3}(u)P(u),\quad B(u)=K(u)Q(u).
\end{align}
Write $a_r=\coef{r}{A}$, $b_r=\coef{r}{B}$, and, for any
polynomials $F,G$ with coefficients $f_i,g_i$, define
\begin{equation}\label{eq:minor}
 \M_r(F,G)=f_rg_{r+1}-f_{r+1}g_r.
\end{equation}

\begin{theorem}[Certified finite-$k$, infinite-$n$ range]\label{thm:main}
For every integer $3\le k\le180$, every integer $n\ge2k+4$,
and every $0\le r\le k$,
\begin{equation}\boxed{\M_r(A,B)>0.}\end{equation}
The upper bound $180$ is a verified computational limit, not a natural
mathematical threshold.
\end{theorem}

\section{Mathematical lemmas}
The Fibonacci matching polynomials satisfy
$Z_m=Z_{m-1}+uZ_{m-2}$ with $Z_0=Z_1=1$. In particular,
\begin{equation}\label{eq:cassini}
Z_{k+1}(u)Z_{k+3}(u)-Z_{k+2}(u)^2=(-1)^{k+1}u^{k+2}.
\end{equation}
Since the correction in \eqref{eq:cassini} has degree at least $k+2$,
\begin{equation}\label{eq:low}
 \M_r(A,B)=\M_r(KP,KQ) \qquad (0\le r\le k).
\end{equation}
The real-negative-root factorization of the matching polynomial reads
\begin{equation}\label{eq:factor}
 Z_m(u)=\prod_{j=1}^{\lfloor m/2\rfloor}
 \left(1+4\cos^2\!\frac{\pi j}{m+1}\,u\right).
\end{equation}
Thus $K(u)=(1+\lambda u)H_0(u)$ with
$\lambda=4\cos^2(\pi/(k+3))\ge3$ for $k\ge3$, where $H_0$
has nonnegative coefficients, positive constant term, and a
P\'olya-frequency coefficient sequence (indeed $\mathrm{PF}_\infty$),
by the classical real-rootedness criterion~\cite{ASW}.

\begin{lemma}[Ratio propagation]\label{lem:ratio}
If $f_i,g_i>0$ for $0\le i\le R+1$ and
$\M_i(F,G)>0$ for all $0\le i\le R$, then
$f_i g_j-f_j g_i>0$ for all $0\le i<j\le R+1$.
\end{lemma}
\begin{proof}
Divide $\M_i(F,G)>0$ by $f_if_{i+1}$ to obtain
$g_i/f_i<g_{i+1}/f_{i+1}$, and chain these inequalities.
\end{proof}

\begin{lemma}[Positive-minor transfer under a common $\mathrm{PF}_2$ convolution]
\label{lem:transfer}
Suppose the hypotheses of Lemma~\ref{lem:ratio} hold. Let
$W(u)=\sum_jw_ju^j$ with $w_0>0$ and nonnegative $\mathrm{PF}_2$
coefficient sequence. Then
$\M_r(WF,WG)>0$ for all $0\le r\le R$.
\end{lemma}
\begin{proof}
The Cauchy--Binet identity for the two Toeplitz-convolved
coefficient columns gives
\begin{align*}
 \M_r(WF,WG)=\sum_{0\le i<j\le r+1}
 &(f_ig_j-f_jg_i)\\[-1mm]
 &\cdot\left(w_{r-i}w_{r+1-j}-w_{r-j}w_{r+1-i}\right),
\end{align*}
where $w_i=0$ outside the support. The first factor is strictly positive;
the second is nonnegative by $\mathrm{PF}_2$. The term
$(i,j)=(r,r+1)$ is $\M_r(F,G)w_0^2>0$.
\end{proof}

\begin{lemma}[One-factor conditional reduction]\label{lem:onefactor}
For any $k\ge3$ and $n\ge2k+4$, assume
\begin{align}
&\M_i(P,Q)>0\quad(0\le i<k), \tag{H1}\label{eq:H1}\\
&\M_k\big((1+\lambda u)P,(1+\lambda u)Q\big)>0,\tag{H2}\label{eq:H2}
\end{align}
where $\lambda=4\cos^2(\pi/(k+3))$. Then the conclusion of
Theorem~\ref{thm:main} holds for this pair $k,n$.
\end{lemma}
\begin{proof}
Apply Lemma~\ref{lem:transfer} first with $W=1+\lambda u$
to the minors of indices $0,\ldots,k-1$ in (H1). The remaining
boundary minor of index $k$ is (H2). Indeed, for $\ell\in\{k-1,k,k+1\}$, every coefficient of
$p_{n,\ell}(u)$ within its support is positive when $n\ge2k+4$.
Consequently,
$\deg P=2\lfloor(k+2)/2\rfloor$ and
$\deg Q=\lfloor(k+1)/2\rfloor+\lfloor(k+3)/2\rfloor$;
each is either $k+1$ or $k+2$. Their coefficients are therefore
positive at every degree $0,\ldots,k+1$, as are those of the
once-smoothed pair. Apply
Lemma~\ref{lem:transfer} a second time with $W=H_0$ and use
\eqref{eq:low}.
\end{proof}

Let $P_i=\coef{i}{P}$ and $Q_i=\coef{i}{Q}$, and define
\begin{align*}
D&=P_kQ_{k+1}-P_{k+1}Q_k,\\
X&=P_{k-1}Q_{k+1}-P_{k+1}Q_{k-1},\\
H&=P_{k-1}Q_k-P_kQ_{k-1}.
\end{align*}
Direct multiplication gives
\begin{equation}\label{eq:J}
J_k(\lambda):=\M_k((1+\lambda u)P,(1+\lambda u)Q)
=D+\lambda X+\lambda^2H.
\end{equation}
Consequently
\begin{equation}\label{eq:Jplus}
J_k(1+t)=J_k(1)+J'_k(1)t+Ht^2.
\end{equation}
Since $\lambda\ge1$, it suffices for (H2) to verify
$H>0$, $J_k(1)>0$, and $J'_k(1)\ge0$.

\section{Exact certificate and proof of Theorem~\ref{thm:main}}
Set $s=n-2k-4\ge0$. For any fixed integer $k$, all expressions in
(H1) and the three sufficient boundary inequalities are polynomials
in $s$ with integer coefficients and degree at most $4k$.
More precisely, the coefficients of $p_{n,k-1}$, $p_{n,k}$ and
$p_{n,k+1}$ are polynomials in $s$ of degrees at most $k-1$, $k$
and $k+1$, respectively: each term is a binomial polynomial with
the corresponding lower index. Hence the coefficients of both
$P=p_{n,k}^2$ and $Q=p_{n,k-1}p_{n,k+1}$ have degree at most $2k$
in $s$; the coefficient minors and the three boundary expressions
have degree at most $4k$.

For any such polynomial $f(s)$ the Newton expansion is the identity
\begin{equation}\label{eq:newton}
 f(s)=\sum_{j=0}^{4k}\Delta^jf(0)\binom{s}{j}
 \quad(s\in\mathbb Z_{\ge0}),
\end{equation}
where $\Delta f(s)=f(s+1)-f(s)$. Thus
$\Delta^jf(0)\ge0$ for all $0\le j\le4k$, with
$\Delta^0f(0)>0$, certifies $f(s)>0$ for every integer $s\ge0$.
A further check of $\Delta^{4k+1}f(0)=0$ detects violations of the
claimed degree bound in the sampled values; the analytic degree bound
is also required to lift finite samples to the identity
\eqref{eq:newton}.

The frozen verifier \texttt{verify\_newton\_v023.py} generates, for every
integer $3\le k\le180$, the complete forward-difference sequence
for the $k$ indices in (H1) and the three boundary quantities
$H$, $J_k(1)$ and $J'_k(1)$. These are exact integer operations.
All tested constant coefficients are positive, all other coefficients
nonnegative, and all $(4k+1)$-st differences vanish. Thus (H1) and
(H2) hold for every $n\ge2k+4$ for each stated $k$. Theorem
\ref{thm:main} follows from Lemma~\ref{lem:onefactor}.

The twelve disjoint proof ranges partition $k=3,\ldots,180$;
reproducibility logs record
\begin{center}
\begin{tabular}{lrrr}\toprule
Criterion & positive & zero & negative\\\midrule
All Newton coefficients & 8,036,344 & 16,821 & 0\\
Total checked & \multicolumn{3}{c}{8,053,165}\\\bottomrule
\end{tabular}
\end{center}
Here $7,857,187$ entries concern the raw-prefix inequalities and
$195,978$ the boundary inequalities. All 12 ranges were also
regenerated on a user-operated Windows computer on 9 November 2026;
the sanitized execution summary and hash of the original local log are
provided in the reproducibility package.

\section{Independent cross-checks and a distinct proof route}\label{sec:crosschecks}
A separate verifier \texttt{verify\_correspondence\_v024.py} compares
$p_{n,k}(1)$ with the recurrence obtained from the rational generating
function of~\cite[Thm.~3.5]{ERZ}: $3741$ exact coefficient cells
for $n\le85$ and $45$ direct permutation enumerations for $n\le8$
were consistent. The analytical generating-function identity is
recorded in the frozen mathematical notes; these finite tests are
supplemental evidence, not a substitute for it.

A second implementation, \texttt{verify\_direct\_newton\_v0241.py},
constructs $A,B$ directly from \eqref{eq:Z}--\eqref{eq:p}, without using
the one-factor reduction. The same degree-$4k$ Newton method
certifies the original minors for $3\le k\le22$ and every integer
$n\ge2k+4$. This independent route checks the conclusion on a
smaller, explicitly stated range; it does not increase the upper bound
$180$ of the main theorem.

The one-factor smoothing is not cosmetic. For $(n,k)=(26,11)$,
$\M_k(P,Q)<0$ whereas the once-smoothed boundary minor
$J_k(1)>0$, so the stronger raw boundary claim is false.

\section{Meaning and limitations}\label{sec:meaning}
Since $A$ and $B$ are products of polynomials whose coefficients are
strictly positive throughout their supports, and each has degree at
least $k+1$, we have $a_i,b_i>0$ for $0\le i\le k+1$.
Theorem~\ref{thm:main} therefore implies
$b_i/a_i<b_{i+1}/a_{i+1}$ for $0\le i\le k$.
Writing $A_{\le k+1}(u)=\sum_{i=0}^{k+1}a_iu^i$ and analogously
for $B$, the truncated ratio has the consequence
\begin{equation}\label{eq:ratioresult}
\frac{d}{du}\left(\frac{A_{\le k+1}(u)}{B_{\le k+1}(u)}\right)<0
\qquad (u>0)
\end{equation}
for the stated $k,n$ range. To justify the derivative, put
$\rho_i=b_i/a_i$ and
$\mathcal{A}(u)=A_{\le k+1}(u)$,
$\mathcal{B}(u)=B_{\le k+1}(u)$.
The sequence $\rho_i$ is strictly increasing, and
$\mathcal{B}/\mathcal{A}$ is its weighted mean with positive
weights $a_i u^i$. More explicitly,
\[
u\frac{d}{du}\frac{\mathcal{B}}{\mathcal{A}}
=\frac{\sum_{0\le i<j\le k+1}(j-i)a_ia_j
(\rho_j-\rho_i)u^{i+j}}{\mathcal{A}(u)^2}>0.
\]
Taking reciprocals proves \eqref{eq:ratioresult}.
This is a consequence of the coefficient-minor inequalities,
not the original conjecture.

Specifically, we do \emph{not} claim:
\begin{enumerate}[nosep]
\item positivity for all integers $k\ge3$;
\item the log-concavity of $b_{n,k}(123,132)$ in $k$;
\item a global sign for $A(u)-B(u)$;
\item established novelty or human peer review.
\end{enumerate}
The search for prior instances of this precise finite-order
minor statement is incomplete. The supporting technology of matching
polynomials, P\'olya-frequency sequences, Newton interpolation and
total positivity is classical, and should not be presented as novel.

\section*{Reproducibility and provenance}
The reproducibility archive's \texttt{repro} directory contains the unchanged
original Newton verifier, all twelve saved Newton-range outputs,
correspondence checks, and the direct auxiliary certificate verifier.
Checksums and exact scope are given in \texttt{repro/README.txt}.
The original v0.24 ZIP, the v0.24.1 audit supplement and the raw Windows
log are preserved separately as private provenance evidence. Independent
AI mathematical and editorial audits reported PASS within their reviewed
scope. Human peer review has not been performed.
\textcopyright{} 2026 Ryutaro Yonezu. All rights reserved.

\begin{thebibliography}{9}
\bibitem{ERZ} S. Elizalde, J. Rivera Jr., and Y. Zhuang,
\emph{Counting Pattern-Avoiding Permutations by Big Descents},
Annals of Combinatorics (2026),
\href{https://doi.org/10.1007/s00026-026-00805-1}{doi:10.1007/s00026-026-00805-1};
\href{https://arxiv.org/abs/2408.15111}{arXiv:2408.15111}.
\bibitem{ASW} M. Aissen, I. J. Schoenberg, and A. M. Whitney,
On the generating functions of totally positive sequences I,
\emph{Journal d'Analyse Math\'ematique} \textbf{2} (1952), 93--103.
\end{thebibliography}
\end{document}
